\section{密度、\texorpdfstring{$L^p$}{Lp} 空间、Hahn--Banach 与对偶}\label{sec:03-03}

\subsection{Radon--Nikodym 与 Lebesgue 分解}\label{subsec:3-3-1}
\index{Radon--Nikodym 定理}\index{Lebesgue 分解}

在区间上，密度 $h$ 产生测度
\[
  \nu(E)=\int_E h\,\dd\mu .
\]
反过来，若我们只知道集合函数 $\nu$，能否找回这样的 $h$？一个必要条件立刻从公式中读出：
$\mu(E)=0$ 时必须有 $\nu(E)=0$。这个条件排除了相对于 Lebesgue 测度的 Dirac 质量，却没有
说明它是否已经充分。更麻烦的是，即便密度存在，我们也不应期待逐点唯一；改变零集上的值不会改变
任何积分。

\begin{definition}[绝对连续与奇异]\label{def:3-3-1-1}
设 $\mu,\nu$ 是同一可测空间 $(X,\mathcal A)$ 上的正测度。若
\[
  \mu(E)=0\quad\Longrightarrow\quad \nu(E)=0,
  \qquad E\in\mathcal A,
\]
则称 $\nu$ 关于 $\mu$ 绝对连续，记作 $\nu\ll\mu$。若存在 $N\in\mathcal A$ 使
$\mu(N)=0$ 且 $\nu(X\setminus N)=0$，则称 $\nu$ 与 $\mu$ 互相奇异，记作 $\nu\perp\mu$。
\end{definition}

\begin{definition}[Radon--Nikodym 导数]\label{def:3-3-1-2}
若 $\nu\ll\mu$ 且存在可测函数 $h:X\to[0,\infty]$，使
\[
  \nu(E)=\int_Eh\,\dd\mu\qquad(E\in\mathcal A),
\]
则把 $h$ 称为 $\nu$ 关于 $\mu$ 的 Radon--Nikodym 导数，记作
$h=\dd\nu/\dd\mu$。这个记号总按 $\mu$-a.e. 等价理解。
\end{definition}

\begin{definition}[Lebesgue 分解]\label{def:3-3-1-3}
若
\[
  \nu=\nu_a+\nu_s,\qquad \nu_a\ll\mu,\qquad \nu_s\perp\mu,
\]
则称右式为 $\nu$ 关于 $\mu$ 的 Lebesgue 分解。
\end{definition}

\begin{example}[密度与原子]\label{ex:3-3-1-1}
在 $\R$ 上记 Lebesgue 测度为 $m$。对 $0<p<1$，令
\[
  \nu=p\,m|_{[0,1]}+(1-p)\delta_0.
\]
第一项的密度是 $p\1_{[0,1]}$，所以它关于 $m$ 绝对连续；第二项集中在 $m$-零集
$\{0\}$ 上，所以它与 $m$ 奇异。特别地，$\nu$ 不可能整个写成 $h\,m$：否则
$1-p=\nu(\{0\})=\int_{\{0\}}h\,\dd m=0$。这正是 Lebesgue 分解中两类质量同时出现的
最小模型。
\end{example}


奇异质量未必由原子组成。下面的例子把这两个概念彻底分开。

\begin{example}[Cantor 测度]\label{ex:3-3-1-2}
令 $C:\R\to[0,1]$ 为标准 Cantor 函数在 $[0,1]$ 外作常值延拓，并令 $\gamma$ 为第~2.3 节
Stieltjes 构造给出的测度，即
\[
  \gamma((s,t])=C(t)-C(s).
\]
$C$ 从 $0$ 增长到 $1$，所以 $\gamma(\R)=1$。Cantor 集 $K$ 的补集是可数个两两不交的三分
开区间 $(a_j,b_j)$ 之并，且 $C$ 在每个 $[a_j,b_j]$ 上恒定。对固定 $j$，由测度从下连续与
Stieltjes 增量公式，
\[
 \gamma((a_j,b_j))
 =\lim_{n\to\infty}\gamma((a_j,b_j-1/n])
 =\lim_{n\to\infty}\bigl(C(b_j-1/n)-C(a_j)\bigr)=0,
\]
其中只取使 $b_j-1/n>a_j$ 的 $n$。因此可数可加性给出
$\gamma(\R\setminus K)=\sum_j\gamma((a_j,b_j))=0$，而第~2~章的长度计算给出
$m(K)=0$，所以 $\gamma\perp m$。另一方面，Stieltjes 测度在 $x$ 处的原子质量是跳跃
\[
  \gamma(\{x\})=C(x)-C(x-)=0,
\]
因为 $C$ 连续。因此 $\gamma$ 是无原子的奇异概率；“奇异”与“离散”并不是同义词。
\end{example}

Radon--Nikodym 定理的困难不在于猜一个公式，而在于证明候选密度已经吃掉全部质量。下面的证明
只使用非负积分的单调收敛定理以及有符号测度的 Hahn 分解。

历史上，Radon 与 Nikodym 的结果把一维微积分中“增量来自导数”的现象推广到了抽象测度。
现代教材也常借助 Hilbert 空间的表示定理证明有限情形；那条路线需要先建立更多泛函分析。
这里采用纯测度的最大密度法：先解决有限质量问题，再用可数分块补上无限总质量的缺口。

\begin{theorem}[Radon--Nikodym 定理]\label{theorem:3-3-1-1}
设 $\mu,\nu$ 是 $(X,\mathcal A)$ 上的 $\sigma$-有限正测度，且 $\nu\ll\mu$。则存在可测
$h\ge0$ 使 $\nu=h\mu$，而且 $h$ 在 $\mu$-a.e. 意义下唯一。若 $\nu(X)<\infty$，则
$h\in L^1(\mu)$ 且
\[
  \int_Xh\,\dd\mu=\nu(X).
\]
\end{theorem}

\begin{proof}
先设 $\mu(X),\nu(X)<\infty$。考虑
\[
 \mathcal H=\left\{g:X\to[0,\infty]\text{ 可测}:
     \int_Eg\,\dd\mu\le\nu(E)\text{ 对每个 }E\in\mathcal A\right\},
 \qquad
 \alpha=\sup_{g\in\mathcal H}\int_Xg\,\dd\mu.
\]
$0\in\mathcal H$，且每个 $g\in\mathcal H$ 满足 $\int g\le\nu(X)$，故
$0\le\alpha\le\nu(X)<\infty$。若 $g_1,g_2\in\mathcal H$，令
$A=\{g_1\ge g_2\}$。对任意 $E\in\mathcal A$，
\[
 \int_E(g_1\vee g_2)\,\dd\mu
 =\int_{E\cap A}g_1\,\dd\mu+\int_{E\setminus A}g_2\,\dd\mu
 \le\nu(E\cap A)+\nu(E\setminus A)=\nu(E).
\]
所以 $g_1\vee g_2\in\mathcal H$。

对每个 $n$ 选取 $g_n\in\mathcal H$，使
$\int g_n>\alpha-1/n$；这里对可数个非空候选集作了逐项选择。令
$u_n=g_1\vee\cdots\vee g_n$，则 $u_n\in\mathcal H$ 且 $u_n\uparrow h$，其中 $h$ 可测。
对任意 $E$，单调收敛定理给出
\[
 \int_Eh\,\dd\mu=\lim_n\int_Eu_n\,\dd\mu\le\nu(E),
\]
故 $h\in\mathcal H$，从而 $\int h\le\alpha$。另一方面，$u_n\ge g_n$，所以
$\int h\ge\int u_n\ge\int g_n>\alpha-1/n$ 对每个 $n$ 成立。令 $n\to\infty$，得到
$\int h\ge\alpha$，故 $\int h=\alpha$。

定义集合函数
\[
  \rho(E)=\nu(E)-\int_Eh\,\dd\mu.
\]
因为 $h\in\mathcal H$，右侧对每个 $E$ 都是两个有限数之差且非负。若 $(E_j)$ 两两不交，对每个
$N$ 有有限可加等式
\[
 \rho\!\left(\bigcup_{j=1}^NE_j\right)=\sum_{j=1}^N\rho(E_j).
\]
又因为 $\nu$ 与 $h\mu$ 都是总质量有限的测度，从下连续性给出
\[
 \begin{aligned}
 \rho\!\left(\bigcup_{j=1}^{\infty}E_j\right)
 &=\lim_{N\to\infty}
   \left[\nu\!\left(\bigcup_{j=1}^NE_j\right)
       -\int_{\cup_{j=1}^NE_j}h\,\dd\mu\right]\\
 &=\lim_{N\to\infty}\sum_{j=1}^N\rho(E_j)
 =\sum_{j=1}^{\infty}\rho(E_j).
 \end{aligned}
\]
因此 $\rho$ 是有限正测度。若 $\mu(E)=0$，则 $\nu(E)=0$ 且
$\int_Eh\,\dd\mu=0$，所以 $\rho(E)=0$；因此 $\rho\ll\mu$。我们证明 $\rho=0$。若不然，则
$\rho(X)>0$。若 $\mu(X)=0$，则 $\nu\ll\mu$ 给出 $\nu=0$，进而 $\rho=0$，与
$\rho(X)>0$ 矛盾。因此 $\mu(X)>0$。取
$\varepsilon=\rho(X)/(2\mu(X))>0$，则
\[
  (\rho-\varepsilon\mu)(X)>0;
\]
对有限有符号测度 $\sigma=\rho-\varepsilon\mu$ 应用 Hahn 分解，取正集 $A$ 与负集
$X\setminus A$。由于
$\sigma(X\setminus A)\le0$ 而 $\sigma(X)>0$，有
\[
 (\rho-\varepsilon\mu)(A)>0.
\]
特别地 $\rho(A)>0$；由 $\rho\ll\mu$ 可知 $\mu(A)>0$。正集的定义还给出：对每个可测
$B\subset A$，
\[
  \rho(B)\ge\varepsilon\mu(B).
\]
于是对任意 $E\in\mathcal A$，利用 $\rho$ 的正性，
\[
 \rho(E)\ge\rho(E\cap A)\ge\varepsilon\mu(E\cap A).
\]
因此
\[
 \int_E(h+\varepsilon\1_A)\,\dd\mu
 \le\int_Eh\,\dd\mu+\rho(E)=\nu(E),
\]
即 $h+\varepsilon\1_A\in\mathcal H$。可是
$\int(h+\varepsilon\1_A)=\alpha+\varepsilon\mu(A)>\alpha$，与 $\alpha$ 的定义矛盾。
故 $\rho=0$，有限情形得证。

现在设 $\mu,\nu$ 仅为 $\sigma$-有限。分别取可测覆盖 $(A_r)$、$(B_s)$，使
$\mu(A_r)<\infty$、$\nu(B_s)<\infty$。把所有交集 $A_r\cap B_s$ 排成一列 $(C_n)$，再令
\[
 D_1=C_1,
 \qquad D_n=C_n\setminus\bigcup_{k<n}C_k\quad(n\ge2).
\]
删去空块后，$(D_n)$ 两两不交、并为 $X$，而且每块同时具有有限的 $\mu$-质量和 $\nu$-质量。
把两测度限制到 $D_n$；它们都是有限测度，而且
$\nu|_{D_n}\ll\mu|_{D_n}$。有限情形给出可测密度 $h_n$，在 $D_n$ 外把它置零。令
$h=\sum_nh_n\1_{D_n}$。这是可测函数。对任意 $E$，可数可加性和非负积分对不交分块的可加性给出
\[
 \nu(E)=\sum_n\nu(E\cap D_n)
 =\sum_n\int_{E\cap D_n}h_n\,\dd\mu
 =\int_Eh\,\dd\mu.
\]

最后证明唯一性。设 $h,k$ 都是密度。在每个有限 $\nu$-质量分块 $D_j$ 上，两者的积分有限，
故它们取无穷值的集合都是 $\mu$-零集；先删去这些集合的可数并，再令
\[
 E_{j,m}=D_j\cap\{h\ge k+1/m\}.
\]
在 $E_{j,m}$ 上两边积分都等于有限数 $\nu(E_{j,m})$，所以
\[
 0=\int_{E_{j,m}}(h-k)\,\dd\mu\ge m^{-1}\mu(E_{j,m}).
\]
故每个 $E_{j,m}$ 都是 $\mu$-零集。对 $j,m$ 取可数并得到
$\mu(\{h>k\})=0$；交换 $h,k$ 即得 $h=k$ a.e.。若 $\nu$ 有限，在表示式中取 $E=X$
便得到最后一项结论。
\end{proof}

任意两个 $\sigma$-有限测度的和仍然 $\sigma$-有限：分别取有限质量覆盖，交叉后即可得到对两者
都有限的可数覆盖。这个小观察使 RN 定理可以直接产生 Lebesgue 分解。

\begin{theorem}[Lebesgue 分解定理]\label{theorem:3-3-1-2}
若 $\mu,\nu$ 是 $(X,\mathcal A)$ 上的 $\sigma$-有限正测度，则存在唯一的正测度
$\nu_a,\nu_s$ 使
\[
  \nu=\nu_a+\nu_s,
  \qquad \nu_a\ll\mu,
  \qquad \nu_s\perp\mu.
\]
而且 $\nu_a=h\mu$，其中 $h\ge0$。
\end{theorem}

\begin{proof}
令 $\tau=\mu+\nu$。由前述交叉分块，$\tau$ 是 $\sigma$-有限测度；又有
$\mu\ll\tau$ 与 $\nu\ll\tau$。由 Radon--Nikodym 定理，存在 $a,b\ge0$ 使
\[
  \mu=a\tau,
  \qquad \nu=b\tau.
\]
两式相加给出，对每个 $E\in\mathcal A$，
\[
 \int_E(a+b)\,\dd\tau=\mu(E)+\nu(E)=\tau(E)=\int_E1\,\dd\tau.
\]
用 RN 导数唯一性，得到 $a+b=1$ $\tau$-a.e.。在这个共同零集上把 $a,b$ 都改为
零；此后 $a,b$ 处处有限，而且在零集外属于 $[0,1]$。
令 $N=\{a=0\}$。因为 $\mu(N)=\int_Na\,\dd\tau=0$，可定义
\[
  h(x)=
  \begin{cases}
    b(x)/a(x),&x\notin N,\\
    0,&x\in N,
  \end{cases}
  \qquad
  \nu_a(E)=\int_Eh\,\dd\mu,
  \qquad
  \nu_s(E)=\int_{E\cap N}b\,\dd\tau.
\]
在 $X\setminus N$ 上有 $ha=b$，所以
\[
 \nu_a(E)=\int_{E\setminus N}\frac ba\,a\,\dd\tau
 =\int_{E\setminus N}b\,\dd\tau.
\]
从而 $\nu=\nu_a+\nu_s$。按定义 $\nu_a\ll\mu$，而 $\nu_s$ 集中在 $\mu$-零集 $N$ 上，
故存在性成立。

设另有 $\nu=h_1\mu+\sigma_1=h_2\mu+\sigma_2$，其中 $\sigma_i\perp\mu$。取
$\mu$-零集 $N_i$ 使 $\sigma_i(X\setminus N_i)=0$，并令 $N=N_1\cup N_2$。在 $X\setminus N$
上，两条等式给出 $h_1\mu=h_2\mu$；在 $N$ 上两个绝对连续测度都为零。因此
$h_1\mu=h_2\mu$ 在整个 $X$ 上成立，继而由总测度相等得到 $\sigma_1=\sigma_2$。
RN 导数的唯一性还给出 $h_1=h_2$ $\mu$-a.e.。
\end{proof}

对复测度，密度还携带一个相位。下一个定理中的模长结论不是记号约定，而是全变差定义的实质后果。

\begin{theorem}[复测度极分解]\label{theorem:3-3-1-3}
若 $\nu$ 是有限复测度，则存在 $|\nu|$-可测函数 $u$，使
\[
  |u|=1\quad |\nu|\text{-a.e.},
  \qquad \dd\nu=u\,\dd|\nu|.
\]
因此，每个关于 $\nu$ 可积的函数 $f$ 都满足
\[
  \int f\,\dd\nu=\int fu\,\dd|\nu|.
\]
\end{theorem}

\begin{proof}
记 $\lambda=|\nu|$。若 $\lambda(E)=0$，则对每个 $A\subset E$ 有
$|\nu(A)|\le\lambda(A)=0$。因此 $\Re\nu$、$\Im\nu$ 的 Jordan 正负部分都关于 $\lambda$
绝对连续。分别应用 Radon--Nikodym 定理并合并实部、虚部，可得某个
$w\in L^1(\lambda;\C)$，使
\[
  \nu(E)=\int_Ew\,\dd\lambda.
\]
在一个 $\lambda$-零集上把 $w$ 改为零后，可将它视为处处有限的复值可测函数。

我们从全变差的分割定义证明
\begin{equation}\label{eq:3-3-1-variation-density}
  |\nu|(E)=\int_E|w|\,\dd\lambda
  \qquad(E\in\mathcal A).
\end{equation}
对 $E$ 的任意有限可测分割 $(E_j)$，
\[
 \sum_j|\nu(E_j)|
 \le\sum_j\int_{E_j}|w|\,\dd\lambda
 =\int_E|w|\,\dd\lambda,
\]
故左边取上确界后得到“$\le$”。为证反向，令
$v=w/|w|$ 于 $\{w\ne0\}$，并在 $\{w=0\}$ 上令 $v=1$。给定
$0<\delta<\pi/2$，把单位圆分成有限个 Borel 弧，使每段的角直径不超过 $\delta$；取每段中点
$\zeta_j$，并令 $E_j=E\cap\{v\text{ 落在第 }j\text{ 段}\}$。于是
\[
 \begin{aligned}
 |\nu|(E)
 &\ge\sum_j\left|\int_{E_j}w\,\dd\lambda\right|
 \ge\sum_j\Re\left(\overline{\zeta_j}\int_{E_j}w\,\dd\lambda\right)\\
 &\ge \cos\delta\sum_j\int_{E_j}|w|\,\dd\lambda
 =\cos\delta\int_E|w|\,\dd\lambda.
 \end{aligned}
\]
令 $\delta\downarrow0$，即得 \eqref{eq:3-3-1-variation-density}。

可是 $\lambda=|\nu|$，故 \eqref{eq:3-3-1-variation-density} 对每个 $E$ 给出
$\int_E|w|\,\dd\lambda=\int_E1\,\dd\lambda$。取
$E_n^+=\{|w|>1+1/n\}$，则
\[
  0=\int_{E_n^+}(|w|-1)\,\dd\lambda\ge n^{-1}\lambda(E_n^+),
\]
故 $\lambda(E_n^+)=0$。同样，对 $E_n^-=\{|w|<1-1/n\}$ 有
\[
  0=\int_{E_n^-}(1-|w|)\,\dd\lambda\ge n^{-1}\lambda(E_n^-),
\]
故 $\lambda(E_n^-)=0$。再对 $n$ 取可数并，得到
$|w|=1$ $\lambda$-a.e.。令 $u=w$ 即得极分解。最后的积分公式先对简单函数由定义成立，
再由第~3.1 节的复积分逼近推广到所有可积 $f$。
\end{proof}

密度可以逐层相乘，但“约掉密度”需要两个测度拥有相同零集。

\begin{proposition}[RN 链式法则与倒数]\label{proposition:3-3-1-4}
设 $\lambda,\nu,\mu$ 均为 $\sigma$-有限正测度。若
$\lambda\ll\nu\ll\mu$，则
\[
 \frac{\dd\lambda}{\dd\mu}
 =\frac{\dd\lambda}{\dd\nu}\frac{\dd\nu}{\dd\mu}
 \quad\mu\text{-a.e.},
\]
其中乘积在 $\dd\nu/\dd\mu=0$ 处取零。若还满足 $\mu\ll\nu$，则
\[
  \frac{\dd\mu}{\dd\nu}
  =\left(\frac{\dd\nu}{\dd\mu}\right)^{-1}
  \quad\nu\text{-a.e.}
\]
\end{proposition}

\begin{proof}
记 $a=\dd\lambda/\dd\nu$、$b=\dd\nu/\dd\mu$。先固定代表。由 $\nu$ 的 $\sigma$-有限性，可用
$\nu(E_j)<\infty$ 的可数集合覆盖 $X$；在每个 $E_j$ 上
$\int_{E_j}b\,\dd\mu<\infty$，故 $b<\infty$ $\mu$-a.e.。在所得共同 $\mu$-零集上把 $b$ 置零。
再用 $\lambda(F_j)<\infty$ 的可数集合覆盖 $X$；等式
$\int_{F_j}a\,\dd\nu=\lambda(F_j)$ 说明 $a<\infty$ $\nu$-a.e.。在所得共同 $\nu$-零集上把
$a$ 置零。
这些修改不改变相应 RN 导数。于是下列乘积都是普通的非负实数乘积。

对非负简单函数 $s$，密度定义直接给出
$\int s\,\dd\nu=\int sb\,\dd\mu$；用单调简单逼近可将此式推广到所有非负可测 $s$。
因此对任意 $E$，
\[
 \lambda(E)=\int_Ea\,\dd\nu=\int_Eab\,\dd\mu.
\]
RN 导数的唯一性给出第一式。

再设 $\mu\ll\nu$。集合 $Z=\{b=0\}$ 满足 $\nu(Z)=0$，故 $\mu(Z)=0$，于是
$b>0$ $\mu$-a.e.。把 $1/b$ 在 $Z$ 上定义为零。对任意 $E$，再次使用非负函数的密度换算，
\[
 \int_E\frac1b\,\dd\nu
 =\int_{E\cap\{b>0\}}\frac1b b\,\dd\mu
 =\mu(E).
\]
这证明倒数公式。反之，$b>0$ $\mu$-a.e. 时，$\nu(E)=0$ 蕴含
$\int_Eb\,\dd\mu=0$，从而 $\mu(E)=0$；所以该正性恰好对应 $\mu$ 与 $\nu$ 等价。
\end{proof}

\begin{example}[链式密度]\label{ex:3-3-1-3}
在 $\R$ 上令
\[
 \dd\nu=2x\1_{(0,1)}(x)\,\dd x,
 \qquad \dd\lambda=6x^2\1_{(0,1)}(x)\,\dd x.
\]
则 $\dd\nu/\dd x=2x\1_{(0,1)}(x)$，而
$\dd\lambda/\dd\nu=3x$ 于 $(0,1)$；在补集上取零。相乘得到
$6x^2\1_{(0,1)}(x)$，即
\[
  \frac{\dd\lambda}{\dd x}
  =\frac{\dd\lambda}{\dd\nu}\frac{\dd\nu}{\dd x}.
\]
补集上的比值只在 $\nu$-零集上被指定，因此不影响链式法则。另一方面，Lebesgue 测度并不关于
$\nu$ 绝对连续，例如 $(2,3)$ 有正 Lebesgue 测度而 $\nu((2,3))=0$；所以不能在全空间写
$\dd x/\dd\nu=(2x\1_{(0,1)})^{-1}$。这说明倒数公式中的测度等价性不能删去。
\end{example}

\begin{figure}[htbp]
  \centering
  \begin{tikzpicture}[node distance=9mm and 13mm]
    \node[book strong node] (nu) {$\nu$\quad 总质量};
    \node[book node,below left=of nu] (ac) {$h\mu$\quad 绝对连续部分};
    \node[book warning node,below right=of nu] (sing) {$\nu_s$\quad 奇异部分};
    \node[book warm node,below=of ac] (base) {$\mu$\quad 参考测度};
    \node[book node,below=of sing] (null) {$N$\quad $\mu(N)=0$};
    \draw[book accent arrow] (nu)--node[left,book label] {密度} (ac);
    \draw[book arrow] (nu)--node[right,book label] {剩余} (sing);
    \draw[book accent arrow] (base)--node[right,book label] {$\times h$} (ac);
    \draw[book arrow] (sing)--node[right,book label] {集中} (null);
  \end{tikzpicture}
  \caption{Lebesgue 分解把能由参考测度看见的质量写成密度；其余质量只能集中在参考测度的零集上。}
  \label{fig:3-3-1-lebesgue-decomposition}
\end{figure}

RN 导数统一了几个看似不同的公式。若 $Q\ll P$ 是两个总质量为一的测度，令
$L=\dd Q/\dd P$。只要 $Q(B)>0$，限制测度归一化便给出
\[
  \frac{Q(A\cap B)}{Q(B)}
  =\frac{\int_{A\cap B}L\,\dd P}{\int_BL\,\dd P}.
\]
这只是密度公式的代入，分母非零的条件不能省略。在概率记号中，左端写作 $Q(A\mid B)$，
$L$ 称为似然比。

再设 $(\Omega,\mathcal F,\Prob)$ 是总质量为一的测度空间，$X\in L^1(\Prob)$，且
$\mathcal G\subset\mathcal F$ 是子 $\sigma$-代数。实值情形中，对
$A\in\mathcal G$ 定义有限有符号测度
\[
  \eta(A)=\int_AX\,\dd\Prob.
\]
对 $A\in\mathcal G$ 令
$\eta_\pm(A)=\int_A X^\pm\,\dd\Prob$。这是有限正测度，且
$\eta_\pm\ll\Prob|_{\mathcal G}$。分别应用 RN 定理，得到 $\mathcal G$-可测函数
$Y_\pm\ge0$，满足 $\eta_\pm=Y_\pm\Prob|_{\mathcal G}$。由于
$\int Y_\pm\,\dd\Prob=\eta_\pm(\Omega)=\int X^\pm\,\dd\Prob<\infty$，函数
$Y=Y_+-Y_-$ 可积，并且
\[
  \int_AY\,\dd\Prob=\int_AX\,\dd\Prob
  \qquad(A\in\mathcal G).
\]
若另一个 $\mathcal G$-可测可积函数 $Z$ 也满足同一等式，对
$A_n^+=\{Y-Z>1/n\}$ 取积分得
$0=\int_{A_n^+}(Y-Z)\,\dd\Prob\ge n^{-1}\Prob(A_n^+)$。对 $n$ 求并并对
$Z-Y$ 重复这一论证，得 $Y=Z$ a.e.。因此 $Y$ 在 a.e. 意义下唯一，记作
$Y_{X,\mathcal G}$。条件期望理论把这一函数记为
$\E[X\mid\mathcal G]$，并进一步研究其正性和 $L^1$ 压缩性。复值情形可分别处理实部与虚部。
若还要把条件积分表示为给定随机变量取值后的概率核，则需要标准 Borel 空间等额外结构；标量
RN 定理本身不保证这种核表示存在。

一维绝对连续函数也落在同一框架。本段采用如下自足定义：
若对每个 $\varepsilon>0$ 都存在 $\delta>0$，使任意有限族两两不交开区间
$(a_k,b_k)\subset[a,b]$ 在 $\sum_k(b_k-a_k)<\delta$ 时满足
$\sum_k|F(b_k)-F(a_k)|<\varepsilon$，就称 $F:[a,b]\to\R$ 绝对连续；并定义
\[
 \operatorname{Var}(F;[a,b])
 :=\sup_{a=t_0<\cdots<t_N=b}\sum_{j=1}^N|F(t_j)-F(t_{j-1})|.
\]
该上确界有限时称 $F$ 有界变差。第~4~章将从微分的角度进一步研究这两个概念。

若 $F:[a,b]\to\R$ 绝对连续，则 $F$ 有界变差。为核实这一点，
对 $\varepsilon=1$ 取绝对连续性给出的 $\delta>0$，再把 $[a,b]$ 分成有限个长度小于
$\delta$ 的区间，设共有 $R$ 段。对任意分割 $a=t_0<\cdots<t_N=b$ 加入这些固定端点。
由三角不等式，加入端点只会使增量绝对值之和不减，所以只需估计加细后的分割。
若固定小区间为 $[c_{r-1},c_r]$，则落在其中的分割子区间两两不交，长度和不超过
$c_r-c_{r-1}<\delta$，所以对应的函数增量绝对值之和小于 $1$。把 $R$ 段估计相加得
\[
 \sum_{j=1}^N|F(t_j)-F(t_{j-1})|<R.
\]
右侧与原分割无关，故 $F$ 有界变差。令 $V(t)$ 为 $F$ 在 $[a,t]$ 上的全变差。
变差对相邻区间可加：把一个分割限制到两个子区间给出一个方向，把两边任意接近上确界的分割
在公共端点拼接给出反向。所以对 $s<t$，
\[
  V(t)-V(s)=\operatorname{Var}(F;[s,t])\ge |F(t)-F(s)|,
\]
故 $(V+F)/2$ 与 $(V-F)/2$ 都单调不减。绝对连续性还说明 $V$ 连续：若 $t-s<\delta$，
$[s,t]$ 内任意分割的子区间总长度小于 $\delta$，所以 $V(t)-V(s)$ 小于相应的
$\varepsilon$。令 $G_+=(V+F)/2$、$G_-=(V-F)/2$，则二者连续且单调不减。在
$[a,b]$ 外将它们作常值延拓，第~2.3 节的 Stieltjes 构造给出有限测度
$\mu_+$、$\mu_-$。令 $\nu_F=\mu_+-\mu_-$，则对 $a\le s<t\le b$，
\[
 \nu_F((s,t])
 =[G_+(t)-G_+(s)]-[G_-(t)-G_-(s)]
 =F(t)-F(s).
\]
还需核对 $\nu_F\ll m$。测度 $\tau=\mu_++\mu_-$ 支配 $|\nu_F|$，且
\[
 \tau((s,t])=V(t)-V(s)=\operatorname{Var}(F;[s,t]).
\]
给定 $\varepsilon>0$，取绝对连续性对应的 $\delta$。若 $m(E)=0$，取开集
$O\supset E$ 使 $m(O)<\delta$，并把 $O\cap[a,b]$ 写成可数个两两不交的相对开区间
$(I_j)$ 之并，并把 $I_j$ 在 $[a,b]$ 中的两个端点记为 $\alpha_j<\beta_j$。对前 $N$ 个分支，
在每个 $[\alpha_j,\beta_j]$ 内取任意有限分割；所有分割子区间的内部仍两两不交，
且总长度不超过 $m(O)<\delta$。令每个分支内的分割逼近相应变差的上确界，绝对连续性因而给出
\[
 \sum_{j=1}^N\operatorname{Var}(F;[\alpha_j,\beta_j])\le\varepsilon.
\]
由于 $G_\pm$ 连续，它们的 Stieltjes 测度在分支端点没有原子，所以
\[
 \tau(I_j)=\operatorname{Var}(F;[\alpha_j,\beta_j]).
\]
令 $N\to\infty$，再用测度的从下连续性，得
$\tau(O\cap[a,b])=\sum_j\tau(I_j)\le\varepsilon$。由于常值延拓不在端点产生原子，
\[
 |\nu_F|(E)\le\tau(E)\le\tau(O\cap[a,b])\le\varepsilon.
\]
令 $\varepsilon\downarrow0$ 得 $|\nu_F|(E)=0$。于是 $\nu_F$ 的 Jordan 两部分都关于 $m$
绝对连续。分别应用 RN 定理并作差，
得到 $h\in L^1([a,b])$，使
\[
  F(t)-F(s)=\int_s^t h(x)\,\dd x.
\]
反向也成立：若 $h\in L^1$，先取 $M>0$ 使
$\int_{\{|h|>M\}}|h|<\varepsilon/2$，再令
$\delta=\varepsilon/(2M)$。若 $I_1,\dots,I_N\subset[a,b]$ 两两不交且
$\sum_{k=1}^N|I_k|<\delta$，记 $A=\bigcup_{k=1}^NI_k$，则
\[
 \begin{aligned}
 \sum_{k=1}^N\left|\int_{I_k}h\,\dd m\right|
 &\le \int_A|h|\,\dd m\\
 &\le \int_{A\cap\{|h|\le M\}}M\,\dd m
      +\int_{\{|h|>M\}}|h|\,\dd m\\
 &\le M\sum_{k=1}^N|I_k|+\frac{\varepsilon}{2}
 <M\delta+\frac{\varepsilon}{2}=\varepsilon.
 \end{aligned}
\]
因此其不定积分绝对连续。第~4 章的微分定理还将进一步识别 $h=F'$ a.e.

这里的边界值得单独记住：绝对连续是 RN 表示的必要条件；密度只在 a.e. 意义下确定；在无限总质量
情形，证明必须落到同时具有有限 $\mu$、$\nu$ 质量的可数分块上，不能对两个无穷总质量作差。

\begin{exercise}[混合测度的分解]\label{exercise:3-3-1-1}
令 $\nu=\tfrac13m|_{[-1,1]}+2\delta_0+\delta_2$。分别求它关于
$m$ 以及关于 $m+\delta_0$ 的 Lebesgue 分解和 RN 密度。
\end{exercise}

\begin{exercise}[RN 导数的零集]\label{exercise:3-3-1-2}
设 $\nu\ll\mu$ 且 $h=\dd\nu/\dd\mu$。证明
$\nu\perp\mu$ 当且仅当 $h=0$ $\mu$-a.e.；解释为什么在两个测度都非零时，这与
$\nu\ll\mu$ 同时成立会造成矛盾。
\end{exercise}

\begin{exercise}[链式法则]\label{exercise:3-3-1-3}
不直接引用命题~\ref{proposition:3-3-1-4}，从简单函数的密度换算公式出发重证链式法则，并说明
三个导数各自在哪个测度的 a.e. 意义下定义。
\end{exercise}

\begin{exercise}[复测度的模长]\label{exercise:3-3-1-4}
设 $\nu=\sum_{j=1}^Nc_j\delta_{x_j}$，其中 $x_j$ 两两不同。计算 $|\nu|$ 和极分解中的 $u$，
并直接用有限分割定义验证答案。
\end{exercise}

密度函数已经出现；接下来需要一个能够精确衡量它们大小、容纳极限并支持积分配对的函数空间。


\subsection{\texorpdfstring{$L^p$}{Lp} 不等式、完备性与几何}\label{subsec:3-3-2}
\index{$L^p$ 空间}\index{Holder@\Holder{} 不等式}\index{Minkowski 不等式}

点态上很高但支撑很窄的函数，究竟算“大”还是“小”？答案取决于我们给尖峰多少惩罚。对
$\alpha\in\R$、$\beta>0$，先看
\[
  f_n=n^\alpha\1_{(0,n^{-\beta})}\qquad((0,1)\text{ 上}).
\]
若以 $p$ 次平均衡量它，则
\begin{equation}\label{eq:3-3-2-spike-norm}
  \int_0^1|f_n|^p\,\dd x=n^{\alpha p-\beta},
  \qquad
  \|f_n\|_p=n^{\alpha-\beta/p}.
\end{equation}
同一列函数可以在一个指数下趋于零、在另一个指数下发散。积分的代数因此要与指数的几何配合起来。

\begin{definition}[$L^p$ 空间]\label{def:3-3-2-1}
对 $1\le p<\infty$，令 $L^p(X,\mathcal A,\mu)$ 为满足
$\int|f|^p\,\dd\mu<\infty$ 的可测实（或复）函数关于 a.e. 相等的商空间，并定义
\[
  \|f\|_p=\left(\int_X|f|^p\,\dd\mu\right)^{1/p}.
\]
\end{definition}

\begin{definition}[$L^\infty$]\label{def:3-3-2-2}
可测函数 $f$ 的本质上确界为
\[
 \|f\|_\infty
 =\esssup_X|f|
 =\inf\{M\ge0:|f|\le M\ \mu\text{-a.e.}\}.
\]
本质有界函数关于 a.e. 相等的商空间记作 $L^\infty$。
\end{definition}

\begin{definition}[共轭指数]\label{def:3-3-2-3}
若 $p,q\in[1,\infty]$ 且 $1/p+1/q=1$，则称 $p,q$ 互为共轭指数；约定
$1/\infty=0$。
\end{definition}

\begin{definition}[Banach 空间]\label{def:3-3-2-banach}
若范数空间中的每个范数 Cauchy 序列都收敛到该空间中的元素，则称它为 Banach 空间。
\end{definition}

有限维范数空间以及配以上确界范数的 $C(K)$ 都是 Banach 空间。完备性只保证 Cauchy 近似有极限，
并不提供有界序列的收敛子列：$\ell^2$ 中的标准单位向量满足
$\|e_j-e_k\|_2=\sqrt2$（$j\ne k$），所以它们甚至没有 Cauchy 子列。

$L^2$ 的早期发展与 Fourier 系数及 Hilbert 空间紧密相连，随后 Riesz 把这种平方可积几何
推广到 $L^p$ 及其对偶。把 $L^p$ 只描述成 Riemann 可积函数的完备化，会遮住测度密度与积分配对
在其形成中的作用。不等式与完备性给出空间的几何结构，对偶表示则识别其中所有连续线性观测。

\begin{example}[尖峰的 $p$ 敏感性]\label{ex:3-3-2-1}
在 \eqref{eq:3-3-2-spike-norm} 中取 $\alpha=1/3$、$\beta=1$。则
\[
 \|f_n\|_2=n^{-1/6}\longrightarrow0,
 \qquad
 \|f_n\|_3=1,
 \qquad
 \|f_n\|_4=n^{1/12}\longrightarrow\infty.
\]
所以“支撑趋于零”并不单独决定平均大小；振幅与指数共同决定结果。
\end{example}

\begin{lemma}[Young 不等式]\label{lemma:3-3-2-1}
设 $a,b\ge0$，$1<p,q<\infty$ 且 $1/p+1/q=1$，则
\[
  ab\le\frac{a^p}{p}+\frac{b^q}{q}.
\]
\end{lemma}

\begin{proof}
固定 $b\ge0$，考虑 $\Phi(t)=t^p/p-bt$，$t\ge0$。若 $b=0$，所需不等式化为
$0\le a^p/p$。若 $b>0$，
$\Phi'(t)=t^{p-1}-b$，故唯一极小点为
$t_0=b^{1/(p-1)}=b^{q-1}$。由于 $(q-1)p=q$ 且 $q/p=q-1$，
\[
 \Phi(t)\ge\Phi(t_0)
 =\frac{b^q}{p}-b^q=-\frac{b^q}{q}.
\]
令 $t=a$ 并移项即得结论。
\end{proof}

\begin{theorem}[\Holder{} 与 Minkowski]\label{theorem:3-3-2-2}
若 $p,q\in[1,\infty]$ 互为共轭指数，$f\in L^p$、$g\in L^q$，则
\[
  \|fg\|_1\le\|f\|_p\|g\|_q.
  \tag{H}\label{eq:3-3-2-holder}
\]
若 $1\le p\le\infty$ 且 $f,g\in L^p$，则
\[
  \|f+g\|_p\le\|f\|_p+\|g\|_p.
  \tag{M}\label{eq:3-3-2-minkowski}
\]
\end{theorem}

\begin{proof}
先证 \Holder{}。若 $1<p,q<\infty$，且某个范数为零，则相应函数 a.e. 为零，结论成立。
其余情形令 $F=|f|/\|f\|_p$、$G=|g|/\|g\|_q$。逐点应用 Young 不等式并积分，得到
\[
 \int FG\,\dd\mu
 \le\frac1p\int F^p\,\dd\mu+\frac1q\int G^q\,\dd\mu
 =\frac1p+\frac1q=1.
\]
乘回两个范数即得 \eqref{eq:3-3-2-holder}。若 $(p,q)=(1,\infty)$，则
$|fg|\le\|g\|_\infty|f|$ a.e.；积分即可。$(p,q)=(\infty,1)$ 对称。

再证 Minkowski。$p=1$ 时由点态三角不等式积分；$p=\infty$ 时，由
$|f+g|\le\|f\|_\infty+\|g\|_\infty$ a.e. 得到结论。设 $1<p<\infty$，并令
$q=p/(p-1)$。不等式
\[
 |f+g|^p\le2^{p-1}(|f|^p+|g|^p)
\]
先保证 $f+g\in L^p$。若 $\|f+g\|_p=0$，则
$0=\|f+g\|_p\le\|f\|_p+\|g\|_p$；以下设该范数非零。由
\[
 |f+g|^p
 \le |f|\,|f+g|^{p-1}+|g|\,|f+g|^{p-1}
\]
以及已经证明的 \Holder{} 不等式，
\[
 \begin{aligned}
 \|f+g\|_p^p
 &\le(\|f\|_p+\|g\|_p)
       \bigl\||f+g|^{p-1}\bigr\|_q\\
 &=(\|f\|_p+\|g\|_p)\|f+g\|_p^{p-1}.
 \end{aligned}
\]
除以非零的 $\|f+g\|_p^{p-1}$ 得到 \eqref{eq:3-3-2-minkowski}。整个推导只在确认左侧有限后
才使用它，因而没有循环调用三角不等式。
\end{proof}

\begin{theorem}[Riesz--Fischer 完备性]\label{theorem:3-3-2-3}
对任意测度空间以及每个 $1\le p\le\infty$，$L^p$ 都是 Banach 空间。若
$f_n\to f$ 于 $L^p$，则存在子列 $(f_{n_k})$ 使
$f_{n_k}(x)\to f(x)$ 对 $\mu$-a.e. 的 $x$ 成立。
\end{theorem}

\begin{proof}
先设 $1\le p<\infty$，并令 $(f_n)$ 为 $L^p$ 中的 Cauchy 序列。递归选取严格递增的指标
$n_k$，使
\[
  \|f_{n_{k+1}}-f_{n_k}\|_p\le2^{-k}.
\]
记 $g_k=|f_{n_{k+1}}-f_{n_k}|$，$S_N=\sum_{k=1}^Ng_k$。Minkowski 不等式给出
\[
  \|S_N\|_p\le\sum_{k=1}^N2^{-k}\le1.
\]
因为 $S_N\uparrow S:=\sum_{k=1}^\infty g_k$，单调收敛定理给出
\[
  \int S^p\,\dd\mu
  =\lim_N\int S_N^p\,\dd\mu\le1.
\]
所以 $S<\infty$ a.e.，增量级数
\[
 f_{n_1}+\sum_{k=1}^\infty(f_{n_{k+1}}-f_{n_k})
\]
在这些点绝对收敛。它的有限部分和都是可测函数，而收敛点集
$\{S<\infty\}$ 也可测；以其和定义 $f$，并在补集上令 $f=0$，便得到可测函数。若
$H_k=\sum_{j=k}^\infty g_j$。对有限尾和使用 Minkowski，再对其 $p$ 次方用 MCT，得
\[
 \begin{aligned}
 \|H_k\|_p^p
 &=\lim_{N\to\infty}
   \left\|\sum_{j=k}^Ng_j\right\|_p^p\\
 &\le\lim_{N\to\infty}
   \left(\sum_{j=k}^N\|g_j\|_p\right)^p
 \le\left(\sum_{j=k}^{\infty}2^{-j}\right)^p.
 \end{aligned}
\]
因而 $\|H_k\|_p\le\sum_{j=k}^\infty2^{-j}=2^{1-k}$。
又有 $|f-f_{n_k}|\le H_k$ a.e.，故 $f-f_{n_k}\in L^p$ 且
$\|f-f_{n_k}\|_p\le2^{1-k}$。由 $f_{n_k}\in L^p$ 及
$f=f_{n_k}+(f-f_{n_k})$，还得到 $f\in L^p$。原序列为 Cauchy，因此给定 $\varepsilon>0$，先取 $N$，使
$m,n\ge N$ 时 $\|f_n-f_m\|_p<\varepsilon/2$；再选 $k$ 使 $n_k\ge N$ 且
$\|f-f_{n_k}\|_p<\varepsilon/2$。于是对每个 $n\ge N$，三角不等式给出
$\|f_n-f\|_p<\varepsilon$，故 $f_n\to f$ 于 $L^p$。

若 $p=\infty$，同样选取快速子列。对每个 $k$ 选代表后，存在零集 $N_k$，使
\[
 |f_{n_{k+1}}(x)-f_{n_k}(x)|\le2^{-k}\qquad(x\notin N_k).
\]
删去可数并 $N=\bigcup_kN_k$ 后，增量级数由 Weierstrass 判别法在 $X\setminus N$ 上一致收敛；
以其和定义 $f$，并在 $N$ 上令 $f=0$。于是
\[
  |f(x)-f_{n_k}(x)|\le2^{1-k}\qquad(x\notin N).
\]
再取零集 $N_0'$，使
$|f_{n_1}(x)|\le\|f_{n_1}\|_\infty+1$ 于 $X\setminus N_0'$。于是对
$x\notin N\cup N_0'$，
\[
 |f(x)|\le |f_{n_1}(x)|+\sum_{k=1}^\infty2^{-k}
 \le\|f_{n_1}\|_\infty+2.
\]
因此 $f\in L^\infty$，而上面的尾估计还给出
$\|f-f_{n_k}\|_\infty\le2^{1-k}$。给定 $\varepsilon>0$，先由原列的 Cauchy 性取
$N_0$，使 $m,n\ge N_0$ 时 $\|f_n-f_m\|_\infty<\varepsilon/2$；再选 $k$ 使
$n_k\ge N_0$ 且 $\|f-f_{n_k}\|_\infty<\varepsilon/2$。于是对每个 $n\ge N_0$，三角不等式给出
$\|f_n-f\|_\infty<\varepsilon$，
故整列收敛于 $f$。

最后设 $f_n\to f$ 于 $L^p$。当 $p<\infty$ 时选子列使
$\|f_{n_k}-f\|_p\le2^{-k}$。令
$G_N=\sum_{k=1}^N|f_{n_k}-f|$，则 Minkowski 不等式给出
\[
 \|G_N\|_p\le\sum_{k=1}^N2^{-k}\le1.
\]
由 $G_N\uparrow G:=\sum_{k=1}^\infty|f_{n_k}-f|$ 及单调收敛，$G\in L^p$，故
$G(x)<\infty$ a.e.；收敛级数的通项趋于零，即 $f_{n_k}(x)\to f(x)$ a.e.
当 $p=\infty$ 时，也选取 $\|f_{n_k}-f\|_\infty\le2^{-k}$。对每个 $k$ 存在零集 $Z_k$，使
$|f_{n_k}(x)-f(x)|\le2^{-k}$ 于 $X\setminus Z_k$；在
$X\setminus\bigcup_kZ_k$ 上直接令 $k\to\infty$，同样得到逐点收敛。
\end{proof}

完备性与依测度收敛是两件不同的事；上面的“快速子列”把范数信息转成可和的逐点增量，正是二者
之间可重复使用的桥梁。

\begin{proposition}[Jensen 不等式]\label{proposition:3-3-2-4}
设 $\mu(X)=1$，$I\subset\R$ 是区间，$f$ 是取值于 $I$ 的可积实函数，
$\varphi:I\to\R$ 为凸函数，且 $\varphi\circ f$ 可积。则
\[
  \varphi\left(\int_Xf\,\dd\mu\right)
  \le\int_X\varphi(f)\,\dd\mu.
\]
\end{proposition}

\begin{proof}
记 $m=\int f\,\dd\mu$。先说明 $m\in I$。若 $a=\inf I$ 是有限左端点，则
$f\ge a$ a.e.，所以 $m\ge a$；如果 $a\notin I$ 而 $m=a$，则 $f-a>0$ a.e.，零积分判据却会从
$\int(f-a)=0$ 推出 $f=a$ a.e.，矛盾。若 $b=\sup I$ 是有限右端点，则
$f\le b$ a.e. 给出 $m\le b$；如果 $b\notin I$ 而 $m=b$，则 $b-f>0$ a.e.，但
$\int(b-f)=0$ 同样违反零积分判据。无穷端点不可能等于有限实数 $m$。
因此 $m$ 要么是 $I$ 的内点，要么是属于 $I$ 的有限端点。先设 $m$ 是内点。凸函数在内点有支撑直线：存在
$c\in\R$ 使
\[
  \varphi(t)\ge\varphi(m)+c(t-m),
  \qquad t\in I.
\]
具体地，可取任意
\[
  c\in\left[
  \sup_{t<m}\frac{\varphi(m)-\varphi(t)}{m-t},
  \inf_{t>m}\frac{\varphi(t)-\varphi(m)}{t-m}
  \right];
\]
凸性保证左端不超过右端。令 $t=f(x)$ 并积分，得到
\[
 \int\varphi(f)\,\dd\mu
 \ge\varphi(m)+c\int(f-m)\,\dd\mu=\varphi(m).
\]

若 $m$ 是 $I$ 中的左端点 $a$，则 $f-a\ge0$ 且积分为零，故 $f=a$ a.e.，等式成立。若 $m$ 是
$I$ 中的右端点 $b$，则 $b-f\ge0$ 且积分为零，故 $f=b$ a.e.，等式也成立。
\end{proof}

两点的总质量为一的测度空间给出最直接的检验：令
$X=\{0,1\}$、$\mu=\theta\delta_0+(1-\theta)\delta_1$，并取
$f(0)=a$、$f(1)=b$。Jensen 正好
化为
\[
 \varphi(\theta a+(1-\theta)b)
 \le\theta\varphi(a)+(1-\theta)\varphi(b),
\]
即凸性的定义。

\begin{example}[有限测度嵌入]\label{ex:3-3-2-2}
若 $\mu(X)<\infty$ 且 $q>p$，则 $L^q\subset L^p$。当 $q<\infty$ 时，对
$|f|^p$ 与 $1$ 使用指数 $q/p$、$q/(q-p)$ 的 \Holder{} 不等式，得到
\[
 \int|f|^p\,\dd\mu
 \le\left(\int|f|^q\,\dd\mu\right)^{p/q}
      \mu(X)^{1-p/q}.
\]
开 $p$ 次方即
\[
  \|f\|_p\le\mu(X)^{1/p-1/q}\|f\|_q.
\]
$q=\infty$ 时由 $|f|\le\|f\|_\infty$ a.e. 直接得到同一公式。

无限测度空间上两个方向都可能失败。设 $q>p$ 且 $q<\infty$。在 $(0,\infty)$ 上取
\[
 f(x)=x^{-a}\1_{(1,\infty)}(x),
 \qquad \frac1q<a\le\frac1p.
\]
因为 $\int_1^\infty x^{-ar}\,\dd x<\infty$ 当且仅当 $ar>1$，所以
$f\in L^q\setminus L^p$。另一方面，取
\[
 g(x)=x^{-b}\1_{(0,1)}(x),
 \qquad \frac1q\le b<\frac1p.
\]
由 $\int_0^1x^{-br}\,\dd x<\infty$ 当且仅当 $br<1$，可得
$g\in L^p\setminus L^q$。若 $q=\infty$，取 $0<a\le1/p$ 的第一个尾部函数，它属于
$L^\infty\setminus L^p$；再取 $0<b<1/p$ 的第二个原点函数，它属于
$L^p\setminus L^\infty$。
\end{example}

\begin{example}[本质上确界]\label{ex:3-3-2-3}
在 $[0,1]$ 上定义
\[
 f(x)=1\quad(x\ne0),
 \qquad f(0)=10^{100}.
\]
则 $\|f\|_\infty=1$，因为异常值只落在零测单点上。相反，在计数测度下单点不是零集，同一修改
会把 $\ell^\infty$ 范数改成 $10^{100}$。所以 $L^\infty$ 的元素是等价类；未经选择代表，点值
$f(x)$ 并不是空间元素所固有的数据。
\end{example}

当 $0<p<1$ 时，公式 $(\int|f|^p)^{1/p}$ 一般不满足三角不等式。若 $A,B$ 是两个质量为 $1$ 的
不交集合，则
\[
 \|\1_A+\1_B\|_p=2^{1/p}>2=\|\1_A\|_p+\|\1_B\|_p.
\]
这个例子的单位球甚至不是凸集：$\1_A$、$\1_B$ 在单位球中，而它们的中点范数为
$2^{1/p-1}>1$。因此在一般测度空间上，这一范围只能得到拟范数结构；单一原子之类的
退化有限维情形可能例外。以下 $L^p$ 理论均取 $1\le p\le\infty$。

\begin{figure}[htbp]
  \centering
  \begin{tikzpicture}[scale=1.1]
    \draw[->,BookInk!70] (-2.1,0)--(2.1,0) node[right] {$x$};
    \draw[->,BookInk!70] (0,-2.1)--(0,2.1) node[above] {$y$};
    \draw[BookWarning,semithick] (-1.65,-1.65) rectangle (1.65,1.65);
    \draw[BookAccent,semithick] (0,0) circle (1.65);
    \draw[BookWarm,semithick] (0,1.65)--(1.65,0)--(0,-1.65)--(-1.65,0)--cycle;
    \node[book label,text=BookWarm] at (1.08,0.72) {$p=1$};
    \node[book label,text=BookAccent] at (1.55,1.02) {$p=2$};
    \node[book label,text=BookWarning] at (1.72,-1.83) {$p=\infty$};
  \end{tikzpicture}
  \caption{二维 $\ell^p$ 单位球的截面。指数越大，坐标尖峰受到的相对惩罚越弱；有限维单位球的
  形状本身不蕴含无限维空间中的紧性。}
  \label{fig:3-3-2-lp-balls}
\end{figure}

这些结论很快会在几个方向中反复出现。$L^2$ 的完备性使平方可积函数成为 Hilbert 空间，并为
Fourier 的 Plancherel 理论提供容纳极限的空间；\Holder{} 不等式控制参数积分、卷积和两个密度的
配对；Riesz--Fischer 的快速子列方法还可用于证明 Banach 值 $L^p$ 空间的完备性。

\begin{figure}[htbp]
  \centering
  \begin{tikzpicture}[node distance=8mm and 8mm]
    \node[book strong node] (f1) {$f_{n_1}$};
    \node[book node,right=of f1] (f2) {$f_{n_2}$};
    \node[book node,right=of f2] (f3) {$f_{n_3}$};
    \node[book node,right=of f3] (dots) {$\cdots$};
    \node[book warm node,right=of dots] (f) {$f$};
    \draw[book accent arrow] (f1)--node[above,book label] {$\|\Delta_1\|_p\le2^{-1}$} (f2);
    \draw[book accent arrow] (f2)--node[above,book label] {$\|\Delta_2\|_p\le2^{-2}$} (f3);
    \draw[book accent arrow] (f3)-- (dots);
    \draw[book arrow] (dots)--node[above,book label] {$\sum|\Delta_k|<\infty$ a.e.} (f);
  \end{tikzpicture}
  \caption{快速 Cauchy 子列把范数增量压成可和级数；点态绝对收敛产生候选极限，尾和估计再把它
  拉回 $L^p$ 范数。}
  \label{fig:3-3-2-fast-cauchy}
\end{figure}

\begin{exercise}[\Holder{} 的等号]\label{exercise:3-3-2-1}
设 $1<p,q<\infty$，$f\in L^p$、$g\in L^q$ 均非零。沿 Young 不等式的等号条件证明：
\Holder{} 取等号当且仅当存在常数 $c>0$ 使 $|f|^p=c|g|^q$ a.e.，并另行说明复相位对
$\int fg$ 取模等号的影响。
\end{exercise}

\begin{exercise}[$L^p$ 收敛与依测度收敛]\label{exercise:3-3-2-2}
若 $1\le p<\infty$ 且 $f_n\to f$ 于 $L^p$，证明对每个 $\varepsilon>0$，
\[
 \mu(|f_n-f|>\varepsilon)
 \le\varepsilon^{-p}\|f_n-f\|_p^p\longrightarrow0.
\]
给出一个依测度收敛但不在 $L^p$ 收敛的尖峰例。
\end{exercise}

\begin{exercise}[$L^\infty$ 的完备性]\label{exercise:3-3-2-3}
把 Riesz--Fischer 定理的 $p=\infty$ 证明改写为不抽取快速子列的直接证明：先为 Cauchy 列选
一串误差阈值，再删去一个可数并零集，证明代表在该零集外一致 Cauchy。
\end{exercise}

\begin{exercise}[有限测度中的严格性]\label{exercise:3-3-2-4}
在 $([0,1],m)$ 上，对任意 $1\le p<q<\infty$ 构造
$f\in L^p\setminus L^q$。再证明若 $\mu(X)<\infty$ 且 $|f|$ a.e. 为常数，则有限测度嵌入中的
范数估计取等号。
\end{exercise}

$L^p$ 已经是完备的几何空间。还剩两个识别问题：可否用简单或连续函数逼近其中每个元素？所有连续
线性观测是否都来自积分配对？第一个问题由稠密性回答，第二个问题最终又回到 RN 密度。


\subsection{稠密性、可分性与 \texorpdfstring{$L^p$}{Lp} 对偶}\label{subsec:3-3-3}
\index{Hahn--Banach 定理}\index{$L^p$ 对偶}\index{可分性!$L^p$}

抽象的 a.e. 等价类不适合直接计算。我们希望先在有限分割上的简单函数中完成计算，再由范数极限
进入整个空间。在有拓扑的 Radon 空间上，还希望把测试函数换成 $C_c$。另一方面，一个线性观测
若只在某个子空间上给出，也需要知道它能否在不增大范数的前提下延到全空间。这两条路线分别通向
稠密性和 Hahn--Banach 定理，并在 $L^p$ 对偶表示中会合。

Riesz 的 $L^p$ 对偶理论与矩问题一同发展：线性观测先由有限测试确定，再被一个可积密度统一表示。
这也解释了下面证明的顺序。我们不会先假设抽象对偶已经可表示，而是让泛函在示性函数上生成测度，
随后调用刚刚证明的 RN 定理。

\begin{definition}[连续对偶与对偶范数]\label{def:3-3-3-dual}
范数空间 $E$ 的连续对偶 $E^*$ 是所有连续线性映射
$\Lambda:E\to\C$（实空间中以 $\R$ 为值域）组成的空间，范数为
\[
  \|\Lambda\|=\sup_{\|x\|\le1}|\Lambda(x)|.
\]
线性泛函连续当且仅当存在 $C<\infty$ 使
$|\Lambda(x)|\le C\|x\|$ 对所有 $x$ 成立；满足此式的最小 $C$ 就是 $\|\Lambda\|$。
\end{definition}

\begin{definition}[积分配对]\label{def:3-3-3-1}
对共轭指数 $p,q$，定义
\[
  \langle f,g\rangle=\int_X f\overline g\,\dd\mu
\]
（实值情形去掉共轭）。全书约定内积或配对对第一变量线性。因此
$\Lambda_g(f)=\langle f,g\rangle$ 关于 $f$ 线性，而 \Holder{} 不等式给出
$|\Lambda_g(f)|\le\|f\|_p\|g\|_q$。
\end{definition}

\begin{definition}[测试类稠密]\label{def:3-3-3-2}
若子空间 $D\subset L^p$ 满足：对每个 $f\in L^p$ 与 $\varepsilon>0$，存在 $d\in D$ 使
$\|f-d\|_p<\varepsilon$，则称 $D$ 在 $L^p$ 中稠密。若 $D$ 最初由点值函数组成，这里所说的是
它们经 a.e. 商映射所得的像。
\end{definition}

\begin{definition}[对偶表示映射]\label{def:3-3-3-3}
对 $g\in L^q$，令
\[
  J(g)=\Lambda_g,
  \qquad \Lambda_g(f)=\int f\overline g\,\dd\mu.
\]
\Holder{} 保证 $J:L^q\to(L^p)^*$ 有定义。对偶问题是判断 $J$ 是否等距、是否满射。
\end{definition}

在讨论具体的 $L^p$ 前，先解决一般范数空间中的延拓障碍。给定 $x_0\ne0$，直线
$M=\Span\{x_0\}$ 上的公式
$f(\alpha x_0)=\alpha\|x_0\|$ 范数为 $1$；无限维中却没有可供使用的正交投影把它延出去。
一维扩张可以算清，但把所有一维扩张拼成全空间需要一个真正的选择原理。

另一个更尖锐的问题已经出现在 $\ell^\infty$。在
$c_0+\Span\{\mathbf1\}$ 上取出“常数尾巴”的泛函
$x+a\mathbf1\mapsto a$；由 $x_n\to0$ 可知它有界。若它能延到整个 $\ell^\infty$，就会产生一个
不能由 $\ell^1$ 求和配对表示的观测。下面的定理先保证延拓存在，定理后再完成这个反例。

\begin{theorem}[Hahn--Banach 延拓定理]\label{theorem:3-3-3-hahn-banach}
设 $E$ 为实向量空间，$P:E\to\R$ 为次线性函数，即
\[
 P(x+y)\le P(x)+P(y),
 \qquad P(tx)=tP(x)\quad(t\ge0).
\]
若 $M\subset E$ 为线性子空间，线性泛函 $f:M\to\R$ 满足 $f(m)\le P(m)$，则存在
$E$ 上的线性泛函 $F$ 延拓 $f$，并满足 $F(x)\le P(x)$。特别地，实或复范数空间子空间上的
有界线性泛函都可保持范数延拓到全空间。
\end{theorem}

\begin{proof}
先做一维扩张。设 $x_0\notin M$，希望在 $M+\R x_0$ 上定义
\[
  F_\alpha(m+tx_0)=f(m)+t\alpha.
\]
当 $t=1$ 与 $t=-1$ 时，受 $P$ 控制所要求的上下界分别汇成
\[
 L:=\sup_{m\in M}\{f(m)-P(m-x_0)\}
 \le \alpha\le
 U:=\inf_{n\in M}\{P(n+x_0)-f(n)\}.
\]
先核对两个端点之间确实有实数。在下面的不等式中取 $n=0$ 可得
$f(m)-P(m-x_0)\le P(x_0)$，故 $L<\infty$；取 $m=0$ 可得
$P(n+x_0)-f(n)\ge-P(-x_0)$，故 $U>-\infty$。另一方面，对任意 $m,n\in M$，
\[
 f(m)+f(n)=f(m+n)
 \le P(m+n)
 \le P(m-x_0)+P(n+x_0),
\]
故 $f(m)-P(m-x_0)\le P(n+x_0)-f(n)$；先对 $m$ 取上确界、再对 $n$ 取下确界便得
$L\le U$。因而 $[L,U]$ 含有实数，取 $\alpha\in[L,U]$。若 $t>0$，由上界应用于 $m/t$，
\[
 f(m)+t\alpha
 \le tP(m/t+x_0)=P(m+tx_0).
\]
若 $t<0$，写 $s=-t>0$，由下界应用于 $m/s$，
\[
 f(m)-s\alpha
 \le sP(m/s-x_0)=P(m-sx_0).
\]
$t=0$ 正是原假设。因此 $F_\alpha\le P$。

令 $\mathscr E$ 为所有受 $P$ 控制的线性延拓对 $(N,g)$，其中
$M\subset N\subset E$，并按定义域包含且取值相容排序。任意链的定义域之并仍为线性子空间；链上
泛函相容，所以其并给出该链的上界。Zorn 引理于是给出极大元 $(N_*,g_*)$。若
$N_*\ne E$，选 $x_0\in E\setminus N_*$，刚才的一维步骤会把 $g_*$ 继续延拓，违背极大性。
所以 $N_*=E$。这里使用 Zorn 引理，因而使用了相应的选择公理形式；无限次延拓并不是一次有限递归。

若 $E$ 为实范数空间且 $|f(m)|\le C\|m\|$，取 $P(x)=C\|x\|$。定理给出
$F(x)\le C\|x\|$；把 $x$ 换成 $-x$ 又得 $-F(x)\le C\|x\|$，故
$|F(x)|\le C\|x\|$。取 $C=\|f\|$ 可保持范数。

最后设 $E$ 为复范数空间。把 $g=\Re f$ 看成 $M$ 上的实线性泛函。它满足
$\|g\|=\|f\|$：一个方向来自 $|\Re f(x)|\le|f(x)|$；反向则对每个 $x$ 乘一个模为 $1$ 的
标量，使 $f(x)$ 旋转到非负实轴。由实情形把 $g$ 保范延拓为实线性泛函 $G$，再定义
\[
  F(x)=G(x)-iG(ix).
\]
实线性直接给出 $F(ix)=iF(x)$，故 $F$ 复线性；对 $x\in M$，
$G(ix)=\Re(if(x))=-\Im f(x)$，所以 $F(x)=f(x)$。给定 $x$，取 $|\zeta|=1$ 使
$\zeta F(x)=|F(x)|$。复线性与 $\Re F(y)=G(y)$ 给出
\[
 |F(x)|=\Re F(\zeta x)=G(\zeta x)
 \le\|f\|\,\|x\|.
\]
于是复情形也保持范数。
\end{proof}

\begin{corollary}[一点范数化与子空间分离]\label{corollary:3-3-3-norming}
若 $x_0\ne0$ 属于范数空间 $E$，则存在 $\phi\in E^*$ 使
\[
  \|\phi\|=1,
  \qquad \phi(x_0)=\|x_0\|.
\]
更一般地，子空间上的连续线性泛函可保持范数延拓到 $E$。
\end{corollary}

\begin{proof}
在 $\Span\{x_0\}$ 上定义 $\phi_0(\alpha x_0)=\alpha\|x_0\|$。它的范数为 $1$，由
Hahn--Banach 定理保范延拓即可。第二句就是定理的范数空间版本。
\end{proof}

例如，在 $(\R^2,\|\cdot\|_\infty)$ 中，对 $x_0=(1,-1)$ 取
\[
  \phi(a,b)=\frac{a-b}{2}.
\]
由 $|a-b|/2\le\max\{|a|,|b|\}$ 可知 $\|\phi\|\le1$；在 $(1,-1)$ 处取等，故
$\|\phi\|=1$ 且 $\phi(x_0)=1=\|x_0\|_\infty$。这里支撑超平面可以写成公式，定理保证无限维中
也有同样的范数化观测。

现在可以闭合 $L^\infty$ 对偶的第一个失败现场。

\begin{example}[$L^\infty$ 的额外对偶]\label{ex:3-3-3-3}
在 $\N$ 的计数测度上，$L^\infty=\ell^\infty$。令 $c_0$ 为趋于零的序列空间，
$\mathbf1=(1,1,\ldots)$，并在
$M=c_0+\Span\{\mathbf1\}$ 上定义
\[
  F_0(x+a\mathbf1)=a.
\]
由于 $\mathbf1\notin c_0$，分解唯一。又因为 $x_n\to0$，
\[
 |a|=\lim_n|x_n+a|\le\|x+a\mathbf1\|_\infty,
\]
而 $F_0(\mathbf1)=1$，所以 $\|F_0\|=1$。Hahn--Banach 定理把它保范延拓为
$F\in(\ell^\infty)^*$。若存在 $g\in\ell^1$ 使
$F(x)=\sum_nx_n\overline{g_n}$，则
$0=F(e_n)=\overline{g_n}$ 对每个 $n$ 成立，故 $g=0$；但这与
$F(\mathbf1)=1$ 矛盾。因此 $(L^\infty)^*$ 一般严格大于由 $L^1$ 配对产生的部分。
\end{example}

\begin{example}[阶梯函数]\label{ex:3-3-3-1}
设 $\mu$ 是 $\R$ 上的 Radon 测度，$E$ 是有限测度 Borel 集，$1\le p<\infty$。给定
$\varepsilon>0$，正则性给出紧集 $K$ 与开集 $U$，使
\[
  K\subset E\subset U,
  \qquad \mu(U\setminus K)<\varepsilon^p.
\]
在每个有界区间上，$\mu$ 的原子至多可数：质量至少 $1/r$ 的原子只有有限个，再对 $r$ 取可数并。
因此可为每个 $x\in K$ 选开区间 $(a_x,b_x)$，使
\[
  x\in(a_x,b_x),
  \qquad [a_x,b_x]\subset U,
  \qquad \mu(\{a_x\})=\mu(\{b_x\})=0.
\]
端点恰是局部分布函数的连续点。紧性给出有限子覆盖
$I=\bigcup_{j=1}^N(a_j,b_j)$。函数 $s=\1_I$ 是只有有限个跳点的阶梯函数，而且
$K\subset I\subset U$，故
\[
  \|\1_E-s\|_p^p
  =\mu(E\mathbin{\triangle} I)
  \le\mu(U\setminus K)<\varepsilon^p.
\]
这给出一次带明确误差的阶梯逼近；端点取连续点还保证改变开闭端点不改变其 $L^p$ 等价类。
\end{example}

一般的简单函数稠密只依赖测度，而从简单函数升级到连续函数需要 Radon 正则性和局部 Urysohn
截断。

\begin{theorem}[简单函数与 $C_c$ 稠密]\label{theorem:3-3-3-1}
设 $1\le p<\infty$。
\begin{enumerate}
  \item 在任意测度空间上，具有有限测度支撑的简单函数在 $L^p$ 中稠密；
  \item 若 $X$ 是局部紧 Hausdorff 空间，$\mu$ 是其上的 Radon 测度，则
        $C_c(X)$ 在 $L^p(\mu)$ 中的像稠密。
\end{enumerate}
\end{theorem}

\begin{proof}
先设 $f\in L^p$。令
\[
  E_n=\{1/n\le|f|\le n\},
  \qquad f_n=f\1_{E_n}.
\]
在 $E_n$ 上有 $n^{-p}\1_{E_n}\le |f|^p$，故积分单调性直接给出
$\mu(E_n)\le n^p\|f\|_p^p<\infty$。又有
$|f-f_n|^p\to0$ a.e. 且 $|f-f_n|^p\le|f|^p$，支配收敛定理给出
$f_n\to f$ 于 $L^p$。固定 $n$ 后，$f_n$ 有界且支撑测度有限。把其实部、虚部的有界值域分别
分成网格不超过 $1/k$ 的有限区间，并在每个原像上取相应网格值，得到有限值简单函数 $s_k$，使
\[
 |s_k-f_n|\le\frac{\sqrt2}{k}\1_{E_n}\quad\text{a.e.}
\]
（实值情形去掉 $\sqrt2$）。所以
$\|s_k-f_n\|_p\le\sqrt2\mu(E_n)^{1/p}/k\to0$，第一项得证。

现在设 $X$ 局部紧 Hausdorff，$\mu$ 为 Radon 测度。只需逼近一个有限测度集合 $E$ 的示性函数。
给定 $\varepsilon>0$，由内外正则性取紧集 $K$ 和开集 $U$，使
\[
  K\subset E\subset U,
  \qquad \mu(U\setminus K)<\varepsilon^p.
\]
局部 Urysohn 引理给出 $\phi\in C_c(X,[0,1])$，满足
$\phi=1$ 于 $K$ 且 $\supp\phi\subset U$。于是
\[
  |\1_E-\phi|\le\1_{U\setminus K},
  \qquad
  \|\1_E-\phi\|_p<\varepsilon.
\]
有限线性组合再配合三角不等式，可逼近每个有限测度支撑简单函数；第一项完成对任意
$f\in L^p$ 的逼近。
\end{proof}

特别地，在 $\R$ 上对简单函数的每个有限测度水平集使用例~\ref{ex:3-3-3-1}，再作有限线性组合，
可知有限区间示性函数张成的阶梯函数在任意 Radon 测度的 $L^p(\R)$ 中稠密。

\begin{example}[欧氏空间中的光滑逼近]\label{ex:3-3-3-2}
在 $\R^n$ 的 Lebesgue 测度下，上一定理已经给出 $C_c(\R^n)$ 的稠密性。取非负
$\rho\in C_c^\infty(\R^n)$，满足 $\int\rho=1$，并置
$\rho_\varepsilon(x)=\varepsilon^{-n}\rho(x/\varepsilon)$。预期的第二步是
\[
  f_\varepsilon=\rho_\varepsilon*f\in C_c^\infty,
  \qquad \|f_\varepsilon-f\|_p\longrightarrow0
  \quad(f\in C_c).
\]
这个极限由卷积、平移连续性和近似恒等族推出，完整证明见
定理~\ref{theorem:4-3-1-3}。一般 Radon 空间未必具有平移或光滑结构，所以其中适用的稠密类是
$C_c(X)$，不能无条件替换为 $C_c^\infty(X)$。
\end{example}

第二可数性把连续逼近进一步压缩成可数测试族。

\begin{proposition}[$L^p$ 可分性]\label{proposition:3-3-3-2}
若 $X$ 是第二可数的局部紧 Hausdorff 空间，$\mu$ 是 Radon 测度，且
$1\le p<\infty$，则 $L^p(\mu)$ 可分。
\end{proposition}

\begin{proof}
由第~1.4 节，可取相对紧开集组成的可数基 $\mathcal B$，以及紧耗尽
\[
 K_1\subset\interior{K_2}\subset K_2\subset\interior{K_3}\subset\cdots,
 \qquad \bigcup_nK_n=X.
\]
对每一对 $V,U\in\mathcal B$ 且 $\overline V\subset U$，由局部 Urysohn 引理选取
$\psi_{V,U}\in C_c(X,[0,1])$，满足 $\psi_{V,U}=1$ 于 $\overline V$ 且
$\supp\psi_{V,U}\subset U$。由 $K_n\subset\interior{K_{n+1}}$，对紧集 $K_n$ 与开集
$\interior{K_{n+1}}$ 再用一次局部 Urysohn 引理，选取 $\eta_n\in C_c(X,[0,1])$，使
\[
 \eta_n=1\text{ 于 }K_n,
 \qquad \supp\eta_n\subset\interior{K_{n+1}}\subset K_{n+1}.
\]
令 $\mathcal A_0$ 为这些函数生成的、不另添全空间
常数函数 $1$ 的有理系数代数；因而 $\mathcal A_0\subset C_c(X)$。复情形
允许实部、虚部均为有理数的系数，并加入所有生成元的共轭。于是 $\mathcal A_0$ 可数。它的一致闭包
对实数（复情形对复数）数乘封闭，因为任意标量都可由相应的有理标量逼近；故这个闭包是
Stone--Weierstrass 定理所需的实或复代数。

这些帽函数分离点。事实上，若 $x\ne y$，由 Hausdorff 正则性和基的性质可取
$x\in V$、$\overline V\subset U$ 且 $y\notin U$，于是
$\psi_{V,U}(x)=1$、$\psi_{V,U}(y)=0$。固定 $n$，把 $\mathcal A_0$ 限制到紧空间
$K_{n+1}$；该代数分离点，并且某个更大的 $\eta_m$ 在 $K_{n+1}$ 上恒为 $1$。实或复
Stone--Weierstrass 定理因此说明它在 $C(K_{n+1})$ 中一致稠密。

给定 $f\in C_c(X)$，取 $n$ 使 $\supp f\subset K_n$。对任意 $\delta>0$，可取
$a\in\mathcal A_0$ 使
\[
  \sup_{K_{n+1}}|a-f|<\delta.
\]
函数 $d=\eta_na$ 仍属于 $\mathcal A_0$，支撑包含在 $K_{n+1}$，且因为
$\eta_nf=f$，
\[
 \|d-f\|_p
 \le\delta\,\mu(K_{n+1})^{1/p}.
\]
Radon 测度在紧集上有限，所以可令右侧任意小。这证明可数集 $\mathcal A_0$ 在 $C_c$ 的
$L^p$ 像中稠密；再用定理~\ref{theorem:3-3-3-1}，它在整个 $L^p$ 中稠密。
\end{proof}

第二可数性不能从陈述中无声删除。若 $I$ 是不可数集，给它离散拓扑和计数测度，则该测度是 Radon，
而 $L^p=\ell^p(I)$ 中的单位向量族 $(e_i)_{i\in I}$ 满足
$\|e_i-e_j\|_p=2^{1/p}$。任意可数子集都不可能逼近全部这些向量，故 $\ell^p(I)$ 不可分。

现在进入对偶表示。核心步骤是让泛函作用在示性函数上，从而产生一个测度；RN 定理再把这个测度
还原成函数。$\sigma$-有限分块同时保证示性函数属于 $L^p$，并避免对无限质量集合直接操作。

\begin{theorem}[$L^p$ 对偶]\label{theorem:3-3-3-3}
设 $(X,\mathcal A,\mu)$ 为 $\sigma$-有限测度空间，$1\le p<\infty$，$q$ 为共轭指数。
则每个 $\Lambda\in(L^p)^*$ 都唯一表示为
\[
  \Lambda(f)=\int_Xf\overline g\,\dd\mu,
  \qquad g\in L^q,
\]
而且 $\|\Lambda\|=\|g\|_q$。
\end{theorem}

\begin{proof}
取两两不交的可测分块 $X=\bigsqcup_{j=1}^\infty X_j$，使
$\mu(X_j)<\infty$。对 $E\subset X_j$ 定义
\[
  \nu_j(E)=\Lambda(\1_E).
\]
若 $(E_n)$ 是 $X_j$ 中两两不交的可测集，令 $E=\bigcup_nE_n$，则
\[
 \left\|\1_E-\sum_{n=1}^N\1_{E_n}\right\|_p^p
 =\mu\left(E\setminus\bigcup_{n=1}^NE_n\right)\longrightarrow0
\]
（$p=1$ 也包含在此公式中）。由 $\Lambda$ 的连续性，
$\nu_j(E)=\sum_n\nu_j(E_n)$。还要核对全变差有限。若 $(A_\ell)_{\ell=1}^r$ 是 $X_j$ 的有限可测分割，
对每个 $\ell$ 选取 $|\zeta_\ell|=1$ 使
$\zeta_\ell\nu_j(A_\ell)=|\nu_j(A_\ell)|$ （零值时任取）。则
\[
 \sum_{\ell=1}^r|\nu_j(A_\ell)|
 =\left|\Lambda\!\left(\sum_{\ell=1}^r\zeta_\ell\1_{A_\ell}\right)\right|
 \le \|\Lambda\|\,\mu(X_j)^{1/p}.
\]
对分割取上确界，得 $|\nu_j|(X_j)<\infty$。所以 $\nu_j$ 是有限复测度；
实标量情形取 $\zeta_\ell\in\{-1,1\}$，得到有限有符号测度。
若 $\mu(E)=0$，则 $\1_E$ 在 $L^p$ 中为零，故 $\nu_j(E)=0$，即 $\nu_j\ll\mu|_{X_j}$。
而且对每个可测 $A\subset E$ 同样有 $\nu_j(A)=0$；全变差的分割定义于是给出
$|\nu_j|(E)=0$，所以实部、虚部的 Jordan 正负部分也都关于 $\mu|_{X_j}$ 绝对连续。
把 $\Re\nu_j$、$\Im\nu_j$ 分成 Jordan 正负部并分别应用 RN 定理，得到
$h_j\in L^1(\mu|_{X_j};\C)$，使
\[
  \nu_j(E)=\int_Eh_j\,\dd\mu.
\]
实情形的 $h_j$ 为实值。令 $h=h_j$ 于 $X_j$。于是对只跨有限多个分块的简单函数 $s$，
\begin{equation}\label{eq:3-3-3-simple-representation}
  \Lambda(s)=\int_Xsh\,\dd\mu.
\end{equation}
在单个有限测度块上，任何有界可测函数都可由一致有界的简单函数逐点逼近；左侧在 $L^p$ 中收敛，
右侧由 $|h_j|$ 支配收敛。因此 \eqref{eq:3-3-3-simple-representation} 也适用于有限多个块上
有界支撑的可测函数。

设 $1<p<\infty$。记 $C=\|\Lambda\|$、$F_N=\bigcup_{j\le N}X_j$，并令
\[
 A_{N,r}=F_N\cap\{|h|\le r\},
 \qquad
 \phi_{N,r}=\overline{\sgn h}\,|h|^{q-1}\1_{A_{N,r}},
\]
其中 $\sgn h=h/|h|$ 于 $h\ne0$，零处取 $0$。函数 $\phi_{N,r}$ 有界并支撑于有限测度集合，
故可代入上式。置 $I_{N,r}=\int_{A_{N,r}}|h|^q\,\dd\mu<\infty$。因为
$(q-1)p=q$，
\[
 I_{N,r}=\Lambda(\phi_{N,r})
 \le C\|\phi_{N,r}\|_p
 =C I_{N,r}^{1/p}.
\]
若 $I_{N,r}>0$，除以 $I_{N,r}^{1/p}$ 得 $I_{N,r}^{1/q}\le C$；若它为零，同一结论成立。
令 $N,r\to\infty$ 并用单调收敛定理，得到
\[
  \int_X|h|^q\,\dd\mu\le C^q.
\]
故 $h\in L^q$ 且 $\|h\|_q\le C$。

设 $p=1$。固定 $j$ 与 $\varepsilon>0$，令
$A=X_j\cap\{|h|>C+\varepsilon\}$。函数
$\phi=\overline{\sgn h}\1_A$ 属于 $L^1$，且
\[
  \int_A|h|\,\dd\mu
  =\Lambda(\phi)
  \le C\|\phi\|_1=C\mu(A).
\]
若 $\mu(A)>0$，则
\[
 A=\bigcup_{n=1}^{\infty}
   \bigl\{x\in A:|h(x)|\ge C+\varepsilon+1/n\bigr\};
\]
因而其中某个集合 $A_n$ 有正测度。这就给出
\[
 \int_A|h|\,\dd\mu
 \ge(C+\varepsilon)\mu(A)+n^{-1}\mu(A_n)
 >(C+\varepsilon)\mu(A),
\]
与上式矛盾。因此
$|h|\le C+\varepsilon$ 于 $X_j$ 上 a.e. 成立。对 $j$ 取可数并、再令
$\varepsilon\downarrow0$，得到 $h\in L^\infty$ 且 $\|h\|_\infty\le C$。

两种情形中都令 $g=\overline h$；这是第一变量线性约定所要求的共轭。有限测度支撑简单函数在
$L^p$ 中稠密，故 \eqref{eq:3-3-3-simple-representation} 与连续性给出
\[
  \Lambda(f)=\int f h\,\dd\mu=\int f\overline g\,\dd\mu
  \qquad(f\in L^p).
\]
上面的截断测试已经给出 $\|g\|_q\le\|\Lambda\|$；\Holder{} 不等式给出反向估计
$\|\Lambda\|\le\|g\|_q$，故范数相等。

若 $g_1,g_2$ 给出同一泛函，则 $h=\overline{g_1-g_2}$ 满足
$\int_Eh\,\dd\mu=0$ 对每个包含于某个 $X_j$ 的可测 $E$ 成立。分别取
$E=\{\Re h>1/n\}\cap X_j$、$\{\Re h<-1/n\}\cap X_j$，以及虚部对应的集合，可得
$h=0$ a.e.，故 $g_1$ 与 $g_2$ 作为 $L^q$ 元素相等。
\end{proof}

这里的满射性由 RN 定理给出；Hahn--Banach 定理在一般范数空间中的作用，则是保证有足够多的
连续线性泛函。定理的 $p=1$ 端点确实给出 $(L^1)^*=L^\infty$；但
例~\ref{ex:3-3-3-3} 已经说明把箭头反过来写成 $(L^\infty)^*=L^1$ 会失败。对
$1<p<\infty$，在比 $\sigma$-有限更一般的测度空间上还可通过半有限化处理，或在局部化
（localizable）测度空间上陈述相应版本；这里采用上述无需隐藏附加条件的统一表述。

\begin{proposition}[稠密测试类上的收敛升级]\label{proposition:3-3-3-4}
设 $D\subset L^p$ 是稠密线性子空间，$(\Lambda_n)\subset(L^p)^*$ 满足
$\sup_n\|\Lambda_n\|\le C<\infty$，并且对每个 $d\in D$，数列
$\Lambda_n(d)$ 收敛。则其在 $D$ 上的极限唯一延拓为某个
$\Lambda\in(L^p)^*$，且
$\Lambda_n(f)\to\Lambda(f)$ 对每个 $f\in L^p$ 成立。
\end{proposition}

\begin{proof}
在 $D$ 上定义 $\Lambda_0(d)=\lim_n\Lambda_n(d)$。极限保持线性，并且
\[
 |\Lambda_0(d)|\le C\|d\|_p.
\]
若 $f\in L^p$，取 $d_k\in D$ 使 $d_k\to f$。上式说明
$(\Lambda_0(d_k))$ 是 Cauchy 列，其极限与逼近列无关；定义该极限为 $\Lambda(f)$，便得到
$\Lambda\in(L^p)^*$ 且 $\|\Lambda\|\le C$。

给定 $f$ 与 $\varepsilon>0$。若 $C=0$，则每个 $\Lambda_n$ 的范数为零，因而
$\Lambda_n=0=\Lambda$；以下设 $C>0$。取 $d\in D$ 使
$\|f-d\|_p<\varepsilon/(3C)$，再取 $n$ 足够大使
$|\Lambda_n(d)-\Lambda(d)|<\varepsilon/3$。于是
\[
 \begin{aligned}
 |\Lambda_n(f)-\Lambda(f)|
 &\le|\Lambda_n(f-d)|+|\Lambda_n(d)-\Lambda(d)|+|\Lambda(d-f)|\\
 &\le C\|f-d\|_p+\frac\varepsilon3+C\|f-d\|_p<\varepsilon.
 \end{aligned}
\]
范数的一致有界性正是同时控制首尾两项所必需的条件。
\end{proof}

\begin{figure}[htbp]
  \centering
  \begin{tikzpicture}[node distance=13mm and 32mm]
    \node[book node] (simple) {有限测度支撑\\简单函数};
    \node[book strong node,below=of simple] (cc) {$C_c(X)$};
    \node[book warm node,below=of cc] (smooth) {$C_c^\infty(\R^n)$};
    \node[book node,right=of cc] (lp) {$L^p$};
    \node[book warning node,below=of lp] (dual) {$\Lambda\in(L^p)^*$};
    \node[book strong node,below=of dual] (density) {$g\in L^q$};
    \draw[book accent arrow] (simple)--node[left,book label,align=center] {Radon\\正则性} (cc);
    \draw[book dashed arrow] (cc)--node[left,book label,align=center] {mollifier\\第~4.3.1 小节} (smooth);
    \draw[book accent arrow] (simple.east)-- (lp.west);
    \draw[book accent arrow] (cc.east)-- (lp.west);
    \draw[book dashed arrow] (lp)--node[right,book label] {$f\mapsto\int f\overline g$} (dual);
    \draw[book arrow] (dual)--node[right,book label] {RN} (density);
  \end{tikzpicture}
  \caption{左侧是逐层加强的测试类；虚线所示的光滑层要等卷积建立后才可使用。右侧中，连续线性
  观测先在简单函数上产生测度，再由 RN 定理识别为密度。}
  \label{fig:3-3-3-density-duality}
\end{figure}

对偶语言可用来定义弱收敛与弱星收敛：前者测试所有连续线性泛函，后者测试一个指定的前对偶。
命题~\ref{proposition:3-3-3-4} 给出一条常用原则：先在稠密类上验证收敛，再借助一致范数界延伸到
整个空间。条件期望算子、分布与谱测度的逼近论证都会采用这一原则。

\begin{exercise}[简单函数稠密]\label{exercise:3-3-3-1}
给出另一种证明：先把 $f$ 限制到
$E_n=\{1/n\le |f|\le n\}$，再对实值函数的正负部分分别作二进制层逼近；复值情形还要分别处理
实部与虚部。证明 $\mu(E_n)<\infty$，并证明所得有限值简单函数具有有限测度支撑且在 $L^p$ 中收敛到
$f$。
\end{exercise}

\begin{exercise}[$C_c$ 逼近的误差]\label{exercise:3-3-3-2}
设 $s=\sum_{k=1}^Na_k\1_{E_k}$，其中 $E_k$ 两两不交且测度有限。用 Radon 正则性分别构造
$\phi_k\in C_c$，并给出保证
$\|s-\sum_ka_k\phi_k\|_p<\varepsilon$ 的明确误差分配。
\end{exercise}

\begin{exercise}[$L^1$ 对偶的端点测试]\label{exercise:3-3-3-3}
在 $p=1$ 的对偶证明中，把集合
$\{|h|>C+\varepsilon\}$ 分别截到有限测度块上是必要的。写出若直接取整个集合时
$\overline{\sgn h}\1_A$ 可能不属于 $L^1$ 的例子，并完成分块论证。
\end{exercise}

\begin{exercise}[一致有界不可删]\label{exercise:3-3-3-4}
在 $\ell^1$ 中令 $D$ 为有限支撑序列，并定义
$\Lambda_n(x)=n x_n$。证明 $\Lambda_n(d)\to0$ 对每个 $d\in D$ 成立，但存在
$x\in\ell^1$ 使 $(\Lambda_n(x))$ 不收敛。指出命题~\ref{proposition:3-3-3-4} 的哪一项假设失败。
\end{exercise}

范数收敛给出强控制，却不是分析中唯一的极限模式。下一节将比较 a.e. 收敛、依测度收敛与
$L^p$ 收敛，并找出把依测度收敛升级为积分收敛所需的额外条件。
